\chapter{Iloczyny tensorowe}
\label{ch:lin9-iloczyny-tensorowe}
\index{iloczyn tensorowy@Iloczyn tensorowy}
\index{własność uniwersalna@Własność uniwersalna iloczynu tensorowego}
\index{tensor prosty@Tensor prosty}
\index{iloczyn Kroneckera@Iloczyn Kroneckera macierzy}

\noindent\textbf{Cel rozdziału.}
W Rozdziale~\ref{ch:lin8-wieloliniowe} poznaliśmy
odwzorowania wieloliniowe. Zbudujemy przestrzeń
wektorową, w której można je traktować jako
odwzorowania liniowe. Udowodnimy istnienie
i jednoznaczność iloczynu tensorowego,
wyprowadzimy wzory na jego bazy i wymiar
oraz przedstawimy działanie odwzorowań
liniowych i macierzy Kroneckera.
Następnie powiążemy tę konstrukcję
z przestrzeniami odwzorowań, dualami,
sumami prostymi i ilorazami.

W całym rozdziale \(K\) jest ciałem,
\(0_K\ne1_K\), a wszystkie przestrzenie
wektorowe i odwzorowania liniowe są nad \(K\).
Konstrukcja iloczynu tensorowego działa
bez założenia skończonego wymiaru;
twierdzenia o wymiarze i macierzach
mają takie założenie sformułowane jawnie.
Dopuszczamy wymiar \(0\) i dowolną
charakterystykę ciała.

\section{Konstrukcja i własność uniwersalna}
\label{sec:lin9-konstrukcja}

\begin{definicja}[Iloczyn tensorowy]
\label{def:lin9-tensorowy}
Niech \(V,W\) będą przestrzeniami
wektorowymi nad \(K\).
\emph{Iloczynem tensorowym} nazywamy
parę \((T,\tau)\), w której \(T\) jest
przestrzenią wektorową nad \(K\),
a \(\tau\colon V\times W\to T\) jest
odwzorowaniem dwuliniowym takim, że
dla każdej przestrzeni wektorowej \(Z\)
nad \(K\) i każdego odwzorowania
dwuliniowego \(b\colon V\times W\to Z\)
istnieje dokładnie jedno odwzorowanie
liniowe \(\widetilde b\colon T\to Z\)
spełniające
\[
 \widetilde b(\tau(v,w))=b(v,w)
 \qquad(v\in V,\ w\in W).
\]
Warunek ten nazywamy
\emph{własnością uniwersalną}.
Po ustaleniu realizacji
oznaczamy \(T=V\otimes_KW\),
\(\tau(v,w)=v\otimes w\);
element \(v\otimes w\) jest
\emph{tensorem prostym}.
\end{definicja}

\begin{twierdzenie}[Istnienie i rozpiętość przez tensory proste]
\label{thm:lin9-istnienie}
Dla dowolnych przestrzeni wektorowych
\(V,W\) nad \(K\) istnieje
iloczyn tensorowy \(V\otimes_KW\).
Jest on rozpięty przez tensory proste:
\[
 V\otimes_KW=
 \operatorname{span}
 \{v\otimes w:v\in V,\ w\in W\}.
\]
\end{twierdzenie}
\begin{proof}
Oznaczmy \(X=V\times W\).
Niech \(F=K^{(X)}\) będzie przestrzenią
wszystkich funkcji \(f\colon X\to K\)
o skończonym nośniku.
Jej bazę stanowią funkcje
\([v,w]\), równe \(1_K\) w punkcie
\((v,w)\) i \(0_K\) w innych punktach.
Istotnie każda funkcja o skończonym
nośniku ma jednoznaczne rozwinięcie
względem tych funkcji.

Niech \(N\subset F\) będzie
podprzestrzenią rozpiętą przez
wszystkie elementy czterech postaci
\[
 \begin{aligned}
 &[v+v',w]-[v,w]-[v',w],\\
 &[av,w]-a[v,w],\\
 &[v,w+w']-[v,w]-[v,w'],\\
 &[v,aw]-a[v,w],
 \end{aligned}
\]
gdzie \(v,v'\in V\), \(w,w'\in W\)
oraz \(a\in K\).
Zdefiniujmy
\(T:=F/N\) oraz
\(\tau(v,w):=[v,w]+N\).
Wszystkie cztery generatory \(N\)
przechodzą w ilorazie na zero,
więc \(\tau\) jest liniowe
w każdym argumencie.

Ustalmy dowolne dwuliniowe
\(b\colon V\times W\to Z\).
Ponieważ \([v,w]\) tworzą bazę \(F\),
istnieje dokładnie jedno liniowe
\(L_b\colon F\to Z\) takie, że
\(L_b([v,w])=b(v,w)\).
Dwuliniowość \(b\) oznacza,
że \(L_b\) zeruje każdy generator \(N\);
zatem \(N\subset\ker L_b\).
Na mocy własności uniwersalnej
ilorazu z
Twierdzenia~\ref{thm:lin2-wlasnosc-ilorazu}
powstaje jednoznaczne liniowe
\(\widetilde b\colon F/N\to Z\)
o \(\widetilde b(f+N)=L_b(f)\).
W szczególności
\(\widetilde b(\tau(v,w))=b(v,w)\).

Wektory \([v,w]\) rozpinają \(F\),
więc ich klasy \(\tau(v,w)\)
rozpinają \(F/N\).
Jeśli inne liniowe \(R\colon T\to Z\)
spełnia \(R(\tau(v,w))=b(v,w)\),
to \(R\) i \(\widetilde b\)
są równe na zbiorze generatorów
i wobec tego na całym \(T\).
Dowodzi to własności uniwersalnej
i tezy o rozpiętości skonstruowanego
iloczynu tensorowego.
\end{proof}

\begin{twierdzenie}[Jednoznaczność realizacji]
\label{thm:lin9-jednoznacznosc}
Jeżeli \((T_1,\tau_1)\) i
\((T_2,\tau_2)\) mają własność
uniwersalną dla tych samych \(V,W\),
to istnieje dokładnie jeden
izomorfizm liniowy \(J\colon T_1\to T_2\)
spełniający
\(J(\tau_1(v,w))=\tau_2(v,w)\).
W szczególności każda taka realizacja
jest rozpięta przez obrazy \(\tau_i(v,w)\).
\end{twierdzenie}
\begin{proof}
Własność uniwersalna \(T_1\)
zastosowana do \(\tau_2\) daje
jedyne liniowe \(J\colon T_1\to T_2\)
spełniające \(J\circ\tau_1=\tau_2\).
Analogicznie otrzymujemy
\(H\colon T_2\to T_1\) z
\(H\circ\tau_2=\tau_1\).
Złożenie \(H\circ J\) oraz
\(\operatorname{id}_{T_1}\)
mają to samo złożenie z \(\tau_1\).
Jednoznaczność w definicji daje
\(H\circ J=\operatorname{id}_{T_1}\).
Analogicznie
\(J\circ H=\operatorname{id}_{T_2}\),
więc \(J\) jest izomorfizmem.

Dla konstrukcji z poprzedniego
twierdzenia obrazy \(\tau_1\)
rozpinają \(T_1\). Ich obrazy
przez surjektywne \(J\)
rozpinają \(T_2\); są to
dokładnie obrazy \(\tau_2\).
\end{proof}

\begin{figure}[!htbp]
\centering
\begin{tikzpicture}[>=Latex,node distance=28mm and 32mm]
 \node (p) {\(V\times W\)};
 \node[right=of p] (t) {\(V\otimes_KW\)};
 \node[below=of t] (z) {\(Z\)};
 \draw[->] (p)--node[above]{\(\tau\)}(t);
 \draw[->] (p)--node[below left]{\(b\)}(z);
 \draw[->] (t)--node[right]{\(\widetilde b\)}(z);
\end{tikzpicture}
\caption{Każde odwzorowanie dwuliniowe
\(b\) ma dokładnie jedną linearyzację
\(\widetilde b\) przez iloczyn tensorowy.}
\label{fig:lin9-uniwersalna}
\end{figure}

\section{Rachunek tensorów prostych}
\label{sec:lin9-proste}

\begin{stwierdzenie}[Tożsamości dwuliniowości]
\label{prop:lin9-tozsamosci}
Dla wszystkich \(v,v'\in V\),
\(w,w'\in W\) i \(a\in K\)
zachodzą równości
\[
 \begin{aligned}
 (v+v')\otimes w&=v\otimes w+v'\otimes w,\\
 v\otimes(w+w')&=v\otimes w+v\otimes w',\\
 (av)\otimes w&=a(v\otimes w)=v\otimes(aw),\\
 0_V\otimes w&=0_{V\otimes W}=v\otimes0_W.
 \end{aligned}
\]
Każdy tensor jest skończoną sumą
tensorów prostych, lecz taki zapis
nie musi być jednoznaczny.
\end{stwierdzenie}
\begin{proof}
Pierwsze trzy tożsamości
wynikają z dwuliniowości \(\tau\).
Wstawiając \(0_K\) do jednej z nich,
otrzymujemy równości z wektorem
zerowym.
Rozpiętość wykazano w
Twierdzeniu~\ref{thm:lin9-istnienie}.
Niejednoznaczność wynika już z
\((v+v')\otimes w
=v\otimes w+v'\otimes w\).
\end{proof}

\begin{stwierdzenie}[Zerowanie tensora prostego]
\label{prop:lin9-zerowanie}
Dla dowolnych przestrzeni
\(V,W\) nad \(K\)
oraz wektorów \(v\in V,w\in W\)
mamy
\[
 v\otimes w=0\quad\Longleftrightarrow\quad
 v=0_V\ \text{lub}\ w=0_W.
\]
\end{stwierdzenie}
\begin{proof}
Implikacja w lewo wynika
z poprzedniego stwierdzenia.
Jeżeli \(v\ne0\) i \(w\ne0\),
można uzupełnić każdy z wektorów
do bazy Hamela odpowiedniej przestrzeni.
Przy wymiarze nieskończonym
stosujemy
Twierdzenie~\ref{thm:lin-istnienie-bazy-zorn}.
Przypisując \(v\) i \(w\)
wartość \(1_K\), a pozostałym
elementom wybranych baz \(0_K\),
konstruujemy funkcjonały
\(\varphi\in V^*\), \(\psi\in W^*\)
o \(\varphi(v)=\psi(w)=1_K\).
Odwzorowanie
\(b(x,y)=\varphi(x)\psi(y)\)
jest dwuliniowe.
Z własności uniwersalnej
daje funkcjonał \(\widetilde b\)
na \(V\otimes W\) z
\(\widetilde b(v\otimes w)=1_K\).
Wobec tego \(v\otimes w\ne0\).
\end{proof}

\section{Bazy, wymiary i rząd tensora}
\label{sec:lin9-bazy}

\begin{twierdzenie}[Baza iloczynu tensorowego]
\label{thm:lin9-baza}
Niech \(\dim V=n<\infty\),
\(\dim W=m<\infty\).
Niech \(B=(v_1,\ldots,v_n)\)
i \(C=(w_1,\ldots,w_m)\)
będą uporządkowanymi bazami.
Wtedy rodzina
\[
 (v_i\otimes w_j)_{\substack{
 1\leq i\leq n\\1\leq j\leq m}}
\]
uporządkowana leksykograficznie,
najpierw według \(i\), potem \(j\),
jest bazą \(V\otimes_KW\).
Każdy tensor ma dokładnie jeden zapis
\[
 t=\sum_{i=1}^n\sum_{j=1}^m
   c_{ij}v_i\otimes w_j,\qquad c_{ij}\in K,
\]
a \(\dim(V\otimes_KW)=nm\).
\end{twierdzenie}
\begin{proof}
Zapiszmy
\(v=\sum_i x_iv_i\) i
\(w=\sum_j y_jw_j\).
Dwuliniowość daje
\(v\otimes w=\sum_{i,j}
x_iy_jv_i\otimes w_j\).
Ponieważ tensory proste rozpinają
iloczyn tensorowy, wskazana rodzina
również go rozpina.

Niech \(v^1,\ldots,v^n\)
oraz \(w^1,\ldots,w^m\)
będą bazami dualnymi.
Dla każdej pary \((p,q)\)
forma
\(b_{pq}(v,w)=v^p(v)w^q(w)\)
jest dwuliniowa.
Daje zatem jednoznaczny
funkcjonał \(L_{pq}\) na
\(V\otimes W\) spełniający
\[
 L_{pq}(v_i\otimes w_j)
 =\delta_{pi}\delta_{qj}.
\]
Jeśli
\(\sum_{i,j}c_{ij}v_i\otimes w_j=0\),
to zastosowanie \(L_{pq}\)
daje \(c_{pq}=0_K\).
Wszystkie współczynniki
są więc zerowe, czyli
rodzina jest liniowo niezależna.
Ma \(nm\) elementów.
Jeśli \(n=0\) lub \(m=0\),
wszystkie tensory proste są zerowe,
a pusta rodzina jest bazą
przestrzeni zerowej.
\end{proof}

\begin{przyklad}[Współrzędne tensora prostego]
\label{prz:lin9-wspolrzedne}
Niech \(V=W=\mathbb R^2\)
z bazą standardową \(e_1,e_2\).
Dla \(v=2e_1+3e_2\),
\(w=e_1-e_2\) zachodzi
\[
 \begin{aligned}
 v\otimes w
 &=2e_1\otimes e_1-2e_1\otimes e_2\\
 &\quad+3e_2\otimes e_1-3e_2\otimes e_2.
 \end{aligned}
\]
Macierz współczynników
w bazie tensorowej ma postać
\(\begin{pmatrix}2&-2\\3&-3\end{pmatrix}\).
\end{przyklad}

\begin{twierdzenie}[Minimalna liczba tensorów prostych]
\label{thm:lin9-rzad}
Niech \(V,W\) mają skończone
wymiary \(n,m\), a
\(t=\sum_{i,j}c_{ij}v_i\otimes w_j\)
będzie tensorem o macierzy
współczynników
\(C_t=(c_{ij})\in\operatorname{Mat}_{n\times m}(K)\).
Minimalna liczba \(r\in\mathbb N_0\)
składników prostych w zapisie
\(t=\sum_{\ell=1}^r x_\ell\otimes y_\ell\)
jest równa \(\operatorname{rk}C_t\).
Liczba ta jest niezależna
od wyboru baz.
\end{twierdzenie}
\begin{proof}
Jeśli \(t=\sum_{\ell=1}^r
x_\ell\otimes y_\ell\),
a współrzędne \(x_\ell,y_\ell\)
to kolumny \(a_\ell\in K^n\)
i \(b_\ell\in K^m\), to
\(C_t=\sum_{\ell=1}^r
a_\ell b_\ell^{\mathsf T}\).
Obraz każdej macierzy
\(a_\ell b_\ell^{\mathsf T}\)
zawiera się w przestrzeni
rozpiętej przez \(a_\ell\).
Zatem \(\operatorname{rk}C_t\leq r\).

Odwrotnie, niech
\(s=\operatorname{rk}C_t\).
Wybierzmy bazę
\(a_1,\ldots,a_s\)
przestrzeni kolumn \(C_t\).
Współrzędne kolejnych kolumn
w tej bazie tworzą wiersze
\(b_1^{\mathsf T},\ldots,b_s^{\mathsf T}\);
stąd
\(C_t=\sum_{\ell=1}^s
a_\ell b_\ell^{\mathsf T}\).
Niech \(x_\ell,y_\ell\)
mają współrzędne odpowiednio
\(a_\ell,b_\ell\).
Jednoznaczność współrzędnych
w bazie tensorowej daje
\(t=\sum_{\ell=1}^s x_\ell\otimes y_\ell\).
Dla \(s=0\) jest to suma pusta,
a \(t=0\).
Wykazaliśmy dwie przeciwne
nierówności, więc minimum wynosi
\(\operatorname{rk}C_t\).
Minimalna liczba składników
nie odwołuje się do baz,
więc jest niezmiennikiem tensora.
\end{proof}

\begin{przyklad}[Tensor nieprosty]
\label{prz:lin9-nieprosty}
W \(\mathbb R^2\otimes\mathbb R^2\)
tensor
\[
 t=e_1\otimes e_1+e_2\otimes e_2
\]
ma macierz współczynników \(I_2\)
o rzędzie \(2\).
Nie może więc mieć postaci
\(v\otimes w\) dla żadnych
\(v,w\in\mathbb R^2\).
Nie każda suma tensorów prostych
jest jednym tensorem prostym.
\end{przyklad}

\section{Odwzorowania indukowane}
\label{sec:lin9-mapy}

\begin{twierdzenie}[Iloczyn tensorowy odwzorowań]
\label{thm:lin9-odwzorowania}
Niech \(S\colon V\to V'\)
i \(T\colon W\to W'\)
będą odwzorowaniami liniowymi.
Istnieje dokładnie jedno
odwzorowanie liniowe
\[
 S\otimes T\colon V\otimes_KW
       \longrightarrow V'\otimes_KW'
\]
takie, że
\[
 (S\otimes T)(v\otimes w)
 =S(v)\otimes T(w).
\]
Jeśli \(S_1\colon V'\to V''\),
\(T_1\colon W'\to W''\)
są liniowe, to
\[
 (S_1\otimes T_1)\circ(S\otimes T)
 =(S_1\circ S)\otimes(T_1\circ T).
\]
Ponadto
\(\operatorname{id}_V\otimes
 \operatorname{id}_W
 =\operatorname{id}_{V\otimes W}\).
Jeśli \(S,T\) są izomorfizmami,
to \(S\otimes T\) jest
izomorfizmem z odwrotnością
\(S^{-1}\otimes T^{-1}\).
\end{twierdzenie}
\begin{proof}
Funkcja
\(b(v,w)=S(v)\otimes T(w)\)
jest dwuliniowa, gdyż
\(S,T\) są liniowe,
a kanoniczne odwzorowanie
\(V'\times W'\to V'\otimes W'\)
jest dwuliniowe.
Własność uniwersalna
daje jednoznaczne liniowe
\(S\otimes T\).

Obie strony równości dla złożenia
są liniowe i na każdym
\(v\otimes w\) mają wartość
\(S_1(S(v))\otimes T_1(T(w))\).
Z rozpiętości przez tensory proste
wynika równość odwzorowań.
Podobnie sprawdzamy identyczność.
Jeśli istnieją odwrotności \(S^{-1}\)
i \(T^{-1}\), wzór dla złożenia
pokazuje, że
\(S^{-1}\otimes T^{-1}\)
jest dwustronną odwrotnością
\(S\otimes T\).
\end{proof}

\begin{definicja}[Iloczyn Kroneckera]
\label{def:lin9-kronecker}
Niech \(A=(a_{pi})\)
ma rozmiar \(r\times n\),
a \(B=(b_{qj})\) rozmiar
\(s\times m\).
\emph{Iloczyn Kroneckera}
\(A\otimes B\) jest macierzą
o rozmiarze \((rs)\times(nm)\)
z blokami
\[
 A\otimes B=
 \begin{pmatrix}
 a_{11}B&\cdots&a_{1n}B\\
 \vdots&\ddots&\vdots\\
 a_{r1}B&\cdots&a_{rn}B
 \end{pmatrix}.
\]
Jej element indeksowany
parą wierszy \((p,q)\)
i parą kolumn \((i,j)\)
wynosi \(a_{pi}b_{qj}\).
Pary porządkujemy leksykograficznie.
Definicja obejmuje również
puste rozmiary macierzy.
\end{definicja}

\begin{twierdzenie}[Macierz odwzorowania tensorowego]
\label{thm:lin9-macierz-kronecker}
Niech \(S\colon V\to V'\)
i \(T\colon W\to W'\)
będą liniowe między przestrzeniami
skończenie wymiarowymi.
Jeśli ich macierze w
wybranych uporządkowanych bazach
są równe odpowiednio \(A,B\),
to macierz \(S\otimes T\)
w odpowiednich bazach tensorowych
jest równa \(A\otimes B\).
Dla zgodnych rozmiarów zachodzą
\[
 (A_1A)\otimes(B_1B)
 =(A_1\otimes B_1)(A\otimes B),
 \qquad
 (A\otimes B)^{\mathsf T}
 =A^{\mathsf T}\otimes B^{\mathsf T}.
\]
\end{twierdzenie}
\begin{proof}
Jeżeli
\(S(v_i)=\sum_pa_{pi}v'_p\)
i \(T(w_j)=\sum_qb_{qj}w'_q\),
to
\[
 (S\otimes T)(v_i\otimes w_j)
 =\sum_{p,q}a_{pi}b_{qj}
                      v'_p\otimes w'_q.
\]
Zatem element macierzy
w wierszu \((p,q)\) i kolumnie
\((i,j)\) jest \(a_{pi}b_{qj}\),
zgodnie z definicją iloczynu
Kroneckera. Równość dla
mnożenia wynika z własności
składania odwzorowań
\(S\otimes T\).
Wzór dla transpozycji
wynika także bezpośrednio
z porównania elementów
o indeksach \(((i,j),(p,q))\).
\end{proof}

\begin{przyklad}[Macierz rozmiaru cztery]
\label{prz:lin9-kronecker}
Dla
\[
 A=\begin{pmatrix}1&2\\0&3\end{pmatrix},
 \qquad
 B=\begin{pmatrix}0&1\\-1&0\end{pmatrix}
\]
nad \(\mathbb R\) otrzymujemy
\[
 A\otimes B=
 \begin{pmatrix}
 0&1&0&2\\
 -1&0&-2&0\\
 0&0&0&3\\
 0&0&-3&0
 \end{pmatrix}.
\]
Działanie na
\((e_1+e_2)\otimes e_1\)
daje tensor
\((3e_1+3e_2)\otimes(-e_2)\),
o kolumnie współrzędnych
\((0,-3,0,-3)^{\mathsf T}\).
Bezpośrednie mnożenie
\(A\otimes B\) przez
\((1,0,1,0)^{\mathsf T}\)
daje ten sam wynik.
\end{przyklad}

\begin{twierdzenie}[Rząd iloczynu Kroneckera]
\label{thm:lin9-rzad-kronecker}
Dla dowolnych macierzy
\(A\in\operatorname{Mat}_{r\times n}(K)\)
i \(B\in\operatorname{Mat}_{s\times m}(K)\)
zachodzi
\[
 \operatorname{rk}(A\otimes B)
 =\operatorname{rk}A\,\operatorname{rk}B.
\]
\end{twierdzenie}
\begin{proof}
Niech \(S\colon K^n\to K^r\)
i \(T\colon K^m\to K^s\)
mają macierze \(A,B\).
Oznaczmy
\(a=\operatorname{rk}S\),
\(b=\operatorname{rk}T\).
Wybierzmy w dziedzinie \(S\)
bazę zaczynającą się od
bazy dopełnienia \(\ker S\),
a dalej od bazy \(\ker S\).
Obrazy pierwszych \(a\) wektorów
tworzą bazę \(\operatorname{im}S\);
uzupełnijmy ją do bazy
przeciwdziedziny.
Macierz \(S\) ma w tych bazach
blok \(I_a\) i pozostałe
bloki zerowe. Analogicznie
macierz \(T\) ma blok \(I_b\)
i pozostałe bloki zerowe.

W bazach tensorowych
\(S\otimes T\) wysyła
dokładnie \(ab\) wektorów
bazowych odpowiadających
pierwszym \(a\) i \(b\)
indeksom na różne wektory
bazy przeciwdziedziny,
a pozostałe wektory bazowe
na zero. Jego rząd wynosi
zatem \(ab\).
Rząd macierzy nie zależy
od wyboru baz.
W bazach standardowych
macierz \(S\otimes T\)
jest równa \(A\otimes B\),
więc otrzymujemy tezę.
Rachunek pozostaje poprawny,
jeśli któryś rząd wynosi \(0\).
\end{proof}
